\subsection{RQ1.3: Explanatory Synthesis}
\label{sec:analysis-rq13}

RQ1.3 asks how activation-transfer size and compute distribution together
explain the latency and boundary-crossing behaviour observed for the
decomposition choices in RQ1.1 and RQ1.2. The analytical contribution of
this section is not to add a measurement but to fuse the two preceding
readings into a single, scoped explanatory model and to state precisely
what that model does and does not answer. RQ1.1 gave the joint reading
across four coarse boundaries but could not separate the two dimensions
because activation size moved with stage depth across all four
conditions. RQ1.2 supplied the controlled contrast at constant
activation size that RQ1.1 lacked. The supporting internal-timing study
(\cref{sec:rq13}) supplies the structural mechanism that links both
observations to the architecture of ResNet-18~\cite{resnet}.

\paragraph{A two-dimensional account is the smallest model that fits both stages.}
At the coarse level (\cref{tab:rq11-results}), the boundary-crossing
interval decreases monotonically from \(53.81\)~ms to \(2.22\)~ms as
the activation transferred at the boundary decreases from
\(802{,}816\)~B to \(100{,}352\)~B
(\texttt{rq1\_1\_20260515\_111428/condition\_summaries.csv}). At the
refined level (\cref{tab:rq12-results}), holding activation transfer
constant at \(200{,}704\)~B isolates a \(7.81\)~ms boundary-crossing
gap between \texttt{split\_after\_layer3\_block0} and
\texttt{split\_after\_layer3} that the per-segment decomposition
attributes essentially in full to service-B compute, with a non-compute
residual that is the same for both conditions
(\texttt{rq1\_2\_20260515\_112636/cross\_condition.csv}). Neither
dimension on its own fits both observations: activation transfer is
well-fit to the coarse gradient but cannot account for the
wire-cost-equal asymmetry, and compute placement accounts for the
wire-cost-equal asymmetry but cannot account for the order-of-magnitude
transfer gradient between coarse boundaries. The two-dimensional
account that treats activation transfer and compute placement as
\emph{complementary} dimensions is therefore the smallest model
consistent with both stages, and it is consistent with the joint
framing already established in the split-inference literature for
both activation-driven and compute-driven failure
modes~\cite{neurosurgeon,dnnsplit,deeperthings}.

\paragraph{The structural compute profile of ResNet-18 supplies a plausible mechanism.}
The supporting internal-timing study reports that \texttt{layer4}
accounts for \(31.28\%\) of monolithic forward-pass compute, the
\texttt{layer3.1} block accounts for \(9.92\%\) (\(7.48\)~ms), and
the classification head accounts for \(0.44\%\)
(\texttt{internal\_timing\_resnet18\_20260515\_113614/summary.csv};
\cref{tab:rq13-internal-summary,fig:rq13-stage-compute-distribution}).
Two predictions follow. First, a split after \texttt{layer4} leaves
the downstream service with only the classification head and is
therefore expected to be compute-degenerate; this is consistent with
the \(0.62\%\) downstream-compute share measured in RQ1.1
(\texttt{carry\_forward.json}). Second, the wire-cost-equal contrast
in RQ1.2 moves precisely the \texttt{layer3.1} block across the
boundary; the monolithic cost of that block (\(7.48\)~ms) and the
measured service-B compute gap (\(7.81\)~ms) agree to within roughly
\(4\%\). The agreement should not be read as identity, because the
two values come from separate runs (instrumented monolithic profiling
versus split benchmarking) and the residual difference is within the
range explained by run-to-run variability and instrumentation. It
does, however, supply a structural reason why the two wire-cost-equal
conditions differ by the magnitude they do, consistent with the prior
observation that block-level execution time can vary across candidate
boundaries even when their activation footprints are
similar~\cite{pareto_split}.

\paragraph{Why the carry-forward valley is mid-to-late.}
The two-dimensional account also predicts the shape of the
carry-forward region. The earliest boundary is excluded because the
layer-1 activation expands the byte budget the boundary must carry,
exceeding the raw input size in this run --- the activation-transfer
dimension is dominant on the early side. The latest boundary is
excluded because almost all meaningful computation sits before
\texttt{layer4} in the architecture --- the compute-distribution
dimension is dominant on the late side. The mid-to-late region around
\texttt{layer2}, \texttt{layer3.0}, and \texttt{layer3} is the only
band in which both dimensions are simultaneously tolerable, and it is
the band that the carry-forward rule selects from in both stages
(\texttt{carry\_forward.json} in both frozen runs). RQ1.2 further
shows that this band is operationally flat at the available
measurement resolution: the three refined splits cluster in a
\(0.235\)~ms window whose internal ordering is not separated at the
per-iteration \(95\%\) CI level, and the mean separation between the
rule-selected main and the rejected condition is smaller than the
round-to-round range of the rule-selected main itself
(\texttt{round\_consistency.json}). The two-dimensional account
therefore reframes the RQ1.2 tie-break: it is the rule choosing
legitimately within a genuinely flat valley, not a discovery that
\texttt{split\_after\_layer3\_block0} is meaningfully faster than the
other two refined candidates.

\paragraph{Validation.}
The synthesis rests on the union of the RQ1.1 and RQ1.2 validation
arguments together with the validation of the supporting
internal-timing run. Functional parity holds across every split used
in the synthesis (\texttt{max\_abs\_diff} \(=0.0\) at tolerance
\(1.0\times 10^{-5}\), \texttt{parity\_validation.json} in both
frozen runs), so the \(7.81\)~ms wire-cost-equal gap and the coarse
boundary-crossing gradient are not confounded with numerical drift
across the partition. Round-level consistency remains low in both
stages (per-condition round CV between \(0.002\) and \(0.011\) in
RQ1.1 and between \(0.0017\) and \(0.0070\) in RQ1.2,
\texttt{round\_consistency.json}), which is the reporting practice on
which the cross-stage comparison rests~\cite{benchmarking_sc15}. The
internal-timing run is itself functionally equivalent
(\texttt{validation.json}) with measured instrumentation overhead of
\(0.708\)~ms (\(0.94\%\) of total forward-pass time,
\texttt{instrumentation\_overhead.json}), so the structural
percentages used in the mechanism argument are not materially
distorted by profiling. Three limits must be stated. First, the
synthesis inherits the scope of the underlying runs: single host,
single build, fixed seed, CPU-only batch-one localhost gRPC; the
order-of-magnitude correspondence between architectural cost shares
and measured gaps on this host is not evidence of host-independent
generality. Second, the \(7.48\)~ms versus \(7.81\)~ms agreement is
across two runs of different instrumentation; it bounds the
mechanism, it does not prove it. Third, the boundary-crossing metric
is a composite that bundles service-B compute with serialisation,
transport, and deserialisation (\cref{sec:methods-contract}); the
two-dimensional account works because the wire-cost-equal contrast
isolates the compute component, not because boundary-crossing has
been decomposed into independent transport and compute channels in
the general case.

\paragraph{Evaluation.}
The evaluative reading of RQ1.3 has three facets. First, against the
monolithic baseline, the two-dimensional account does \emph{not}
identify a configuration that is faster than monolithic execution
under RQ1.1 and RQ1.2 conditions --- every split is measurably slower
(\(+2.4\%\) to \(+4.7\%\) in RQ1.1; \(+3.34\%\) to \(+3.64\%\) in
RQ1.2). The contribution of RQ1.3 is therefore explanatory, not
optimisation-oriented. Second, against a single-dimension decision
rule, the joint criterion is strictly stronger: a transfer-only rule
would select \texttt{split\_after\_layer4} and produce a structurally
degenerate microservice, and a compute-only rule has no principled
mechanism for rejecting \texttt{split\_after\_layer1}. Third, against
the carry-forward outcome, the joint criterion ratifies the rule's
selection of the mid-to-late band without ratifying any particular
split inside it: the rule's tie-break is a deterministic choice within
an operationally flat valley, and the defensible carry-forward set is
the pair \texttt{split\_after\_layer3\_block0} (rule-selected main)
plus \texttt{split\_after\_layer2} (reference). RQ1.3 does not
establish that this carry-forward set is host-independent; it
establishes that it is the set the joint criterion identifies under
these conditions.

\paragraph{Answer to RQ1.3.}
Under the controlled local conditions of RQ1.1 and RQ1.2, the latency
and boundary-crossing behaviour of the evaluated splits is well
explained by a two-dimensional model: activation-transfer size
dominates at the coarse level and accounts for the monotonic
boundary-crossing gradient across stage
boundaries~\cite{deeperthings}, and compute placement at the boundary
is independently observable at the refined level and accounts for the
wire-cost-equal asymmetry between \texttt{split\_after\_layer3\_block0}
and \texttt{split\_after\_layer3}~\cite{pareto_split}. The supporting
internal-timing study supplies a plausible structural mechanism that
links both observations to the late-stage compute concentration of
ResNet-18~\cite{resnet}: the same architectural feature that makes
\texttt{layer4} a degenerate split, because almost no computation
remains downstream of it, also makes \texttt{layer3.1} a costly block
to relocate across the boundary, because its monolithic cost
approximates the measured wire-cost-equal gap. The carry-forward
region selected by the rule in both stages --- the mid-to-late band
around \texttt{layer2} to \texttt{layer3}, consistent with the
clustering of suitable boundaries at architecturally meaningful points
reported in the split-inference literature~\cite{dnnsplit} --- is the
band in which both dimensions are simultaneously tolerable, and the
carry-forward outcome \texttt{split\_after\_layer3\_block0} (main)
plus \texttt{split\_after\_layer2} (reference) is the rule's
deterministic selection within that band rather than a separation
discovered by RQ1.2. These claims are scoped to single-host,
single-build, fixed-seed, CPU-only batch-one localhost gRPC; whether
the joint criterion is shaped by infrastructure context is the
empirical question that the later Kubernetes and Azure stages address.
